\section{Introducción}

ROS~2 (Robot Operating System 2) es un middleware de código abierto para el desarrollo de sistemas robóticos~\cite{macenski2022robot}. A diferencia de un sistema operativo convencional, ROS~2 no gestiona hardware directamente. Provee una capa de abstracción que permite a procesos independientes, llamados nodos, comunicarse entre sí mediante tópicos, servicios y acciones, sin conocer los detalles de implementación del otro.

Una característica que suele generar confusión en desarrolladores que se inician en ROS~2 es que el sistema no oculta sus capas internas. Para usarlo correctamente es necesario entender cómo el sistema operativo gestiona procesos, qué son las variables de entorno, y cómo un argumento escrito en la terminal termina disponible como un valor numérico dentro de una clase Python.

Este documento traza ese recorrido de forma exhaustiva. Las primeras dos secciones cubren el entorno del sistema operativo: variables de entorno, el comando \texttt{source} y la naturaleza del ejecutable \texttt{ros2}. La sección sobre el launch file explica los mecanismos de sustitución diferida y la creación de descriptores de nodo. La sección sobre el arranque del nodo describe el ciclo de inicialización, el Parameter Server y las comunicaciones DDS. Las secciones finales cubren el paso de parámetros a clases Python puras y el loop de eventos en tiempo real.

A lo largo del documento se utiliza como hilo conductor el siguiente comando, representativo de los experimentos de control de formación desarrollados en este trabajo:

\begin{lstlisting}[style=bash, caption={Comando de ejemplo utilizado a lo largo del documento.}, label={lst:cmd_ejemplo}]
ros2 launch simulador demo_pva.launch.py \
    n_rays_used:=5 vmax:=0.8
\end{lstlisting}

Este comando lanza una simulación en Gazebo con nodos de control que implementan el Polígono de Velocidades Admisibles (PVA)~\cite{ramirez2001collision}. Cada argumento pasado en la terminal recorrerá las cuatro capas descritas antes de llegar al constructor de la clase \texttt{PVA}.
